%\graphicspath{{content/chapters/3_methodology/figures/}}
%
%\chapter{Methodology}%
%\label{chp:methodology}
%
%
%\section{Image Retrieval}
%
%\subsection{Data Acquisition and Subsetting}
%
%Images were retrieved from ReLAION-5B, distributed as 128 shards (approximately 3.62GB each). Due to storage and computational constraints, processing was restricted to the first four shards, yielding an initial subset of 65,552,809 URL--caption pairs. Shards were selected deterministically in index order (i.e., the first four shards). LAION shards are generated during distributed crawling and do not encode semantic ordering; therefore, shard index does not inherently correspond to content category or structure. However, restricting processing to a subset does not guarantee full representativeness. Consequently, all findings derived in this work apply strictly to the processed subset rather than the entire LAION-5B corpus.
%
%\subsection{Image Source Distribution}
%
%An analysis of valid image origins indicated that 18.79\% of all valid images were derived from five domains:
%\begin{itemize}
%	\item cdn.shopify.com: 3,053,782
%	\item i.pinimg.com: 2,585,519
%	\item i.ebayimg.com: 966,820
%	\item images-na.ssl-images-amazon.com: 886,054
%	\item thumbs.dreamstime.com: 678,516
%\end{itemize}
%
%This concentration suggests strong domain centralisation, with large commercial platforms contributing disproportionately. Such concentration may introduce commercial or e-commerce visual bias, as product imagery and stock photography may be overrepresented relative to smaller independent domains.
%
%\subsection{Image Validity Criteria}
%
%Since ReLAION-5B contains URLs rather than stored image binaries, each URL was resolved and downloaded. An image was considered valid only if all of the following conditions were satisfied:
%
%\paragraph{HTTP Resolution}
%\begin{itemize}
%	\item A GET request completed successfully.
%	\item The server returned an HTTP success status (2xx).
%	\item The request completed within a 10-second timeout.
%	\item A browser-like \texttt{User-Agent} header was used.
%	\item A domain-based \texttt{Referer} header was supplied.
%	\item No network exception (timeout, DNS failure, SSL error, connection error) occurred.
%\end{itemize}
%
%\paragraph{Image Decoding}
%\begin{itemize}
%	\item The response content could be opened using \texttt{PIL.Image.open()}.
%	\item The file was not corrupted or truncated.
%	\item The image successfully executed \texttt{image.load()}.
%\end{itemize}
%
%\paragraph{RGB Conversion}
%\begin{itemize}
%	\item The image was successfully converted to RGB using \texttt{image.convert("RGB")}.
%	\item Non-RGB formats (e.g., RGBA, grayscale) were normalised to RGB.
%	\item Images that failed conversion were discarded.
%\end{itemize}
%
%\paragraph{Batch-Level Validation During Embedding}
%\begin{itemize}
%	\item If embedding computation failed for a batch, recursive binary splitting was applied.
%	\item Images that individually passed forward propagation were retained.
%	\item Failed samples were isolated and excluded.
%\end{itemize}
%
%\subsection{Decay Rate}
%
%From 65,552,809 attempted URLs, 43,491,938 valid images were successfully retrieved. The decay rate was computed as:
%
%\[
%\text{Decay Rate} = 1 - \frac{43,491,938}{65,552,809} \approx 0.33
%\]
%
%Decay includes HTTP failures, decoding failures, and conversion failures. This level of link attrition aligns with prior findings in large-scale web datasets and reflects natural web degradation over time. However, link decay may introduce survivorship bias if certain domains or content types disappear at disproportionate rates.
%
%\subsection{Image Preprocessing}
%
%Images were preprocessed using the canonical pipeline associated with \texttt{openai/clip-vit-large-patch14} via the HuggingFace \texttt{AutoProcessor}. Preprocessing steps included:
%
%\begin{itemize}
%	\item Resize to the model's expected resolution (typically 224$\times$224).
%	\item Preservation of aspect ratio prior to cropping.
%	\item Center crop to square.
%	\item Conversion to tensor representation.
%	\item Pixel value rescaling.
%	\item Channel-wise normalisation using CLIP mean and standard deviation.
%	\item Conversion to \texttt{float32} tensors.
%\end{itemize}
%
%No data augmentation was applied. All preprocessing was deterministic to ensure reproducibility.
%
%\subsection{Embedding Computation}
%
%CLIP embeddings were recalculated for all valid images, as original LAION-5B embeddings are no longer publicly available. Embedding extraction was performed using:
%
%\begin{itemize}
%	\item Model: \texttt{openai/clip-vit-large-patch14}
%	\item Pretraining corpus: OpenAI WebImageText
%	\item Embedding dimensionality: 768
%	\item Patch resolution: 14$\times$14
%\end{itemize}
%
%Images were processed in batches (batch size = 200) and forwarded through the CLIP image encoder using \texttt{model.get\_image\_features(...)} under \texttt{torch.no\_grad()} in evaluation mode. GPU acceleration was utilised when available.
%
%Each embedding vector $\mathbf{e}$ was L2-normalised:
%
%\[
%\mathbf{e}_{norm} = \frac{\mathbf{e}}{\|\mathbf{e}\|_2}
%\]
%
%This ensures unit-length representations and enables cosine similarity search via dot-product equivalence. Embeddings were stored in \texttt{float32} format and written incrementally to Parquet files with checkpointing for fault tolerance and resumability.
%
%\subsection{Model Selection Justification}
%
%While the original LAION-5B filtering pipeline relied on CLIP ViT-B/32 embeddings, this work employs the higher-capacity ViT-L/14 variant for retrieval and bias analysis. ViT-L/14 provides increased representational capacity (768-dimensional embeddings), finer spatial granularity (14$\times$14 patches versus 32$\times$32), and improved semantic alignment in retrieval benchmarks. As the objective of this thesis is not to replicate LAION filtering thresholds but to analyse representational bias, the use of a higher-capacity backbone is methodologically appropriate. Nonetheless, embeddings inherit inductive biases from their pretraining corpus, and retrieval behaviour remains model-dependent.
%
%\subsection{Computational Environment}
%
%The download and embedding pipeline was executed on an NVIDIA RTX 3060 Ti GPU with 64GB system RAM. Processing spanned approximately two to three months. The extended duration may introduce minor temporal drift in URL availability.
%
%\subsection{Ethical Considerations}
%
%\begin{itemize}
%	\item Images were accessed exclusively via publicly available URLs.
%	\item No authentication or access restrictions were bypassed.
%	\item Images were used strictly for academic research purposes.
%	\item No redistribution of image binaries was performed.
%	\item Analysis was conducted at aggregate statistical levels.
%	\item No attempt was made to identify individuals.
%	\item No facial recognition or re-identification techniques were applied.
%	\item Where possible, only embeddings were retained for analytical tasks.
%	\item The research objective is to analyse bias rather than amplify or deploy it.
%\end{itemize}
%
%\subsection{Limitations}
%
%\begin{itemize}
%	\item Only 4 of 128 shards were processed; findings are limited to this subset.
%	\item Shard selection does not guarantee full representativeness.
%	\item Approximately 33\% of URLs were unavailable, introducing survivorship bias.
%	\item 18.79\% of valid images originated from five domains, indicating domain concentration.
%	\item Embeddings reflect biases inherent to CLIP pretraining data.
%	\item Image resizing to 224$\times$224 may reduce fine-grained demographic cues.
%	\item A 10-second timeout may classify slow servers as invalid.
%	\item No geographic replication of downloads was performed.
%	\item Results are encoder-dependent and not model-agnostic.
%	\item No additional content filtering was applied prior to embedding extraction.
%\end{itemize}





%\subsection{Profession-Conditioned Subset Construction}
%
%Following large-scale image retrieval and embedding computation, a structured subset of images was required upon which spurious feature analysis would be conducted. In line with prior work examining occupational representation and stereotype formation in web-scale datasets (Section~\ref{subsec:occupations}), the decision was made to focus on profession-conditioned imagery. A set of 200 diverse professions was constructed and aligned with established ISCO-based occupational categories. For clarity and readability, the complete profession list is provided in the Appendix.
%
%Previously computed CLIP ViT-L/14 image embeddings (768-dimensional, L2-normalised) were used to retrieve profession-aligned images. Two retrieval strategies were evaluated. First, a brute-force similarity computation was implemented, wherein each image embedding was compared against the corresponding CLIP text embedding using cosine similarity. This approach has time complexity $\mathcal{O}(n)$, where $n$ denotes the number of valid images (over 43 million). Although exact, this method proved computationally infeasible at scale.
%
%Consequently, FAISS was adopted (Section~\ref{subsec:faiss}). An \texttt{IndexFlatIP} index was constructed over the full embedding set. Since embeddings were L2-normalised, inner product similarity is equivalent to cosine similarity. While \texttt{IndexFlatIP} performs exact search, FAISS provides highly optimised vectorised implementations, resulting in substantial practical speed improvements relative to naïve brute-force computation.
%
%\subsection{Retrieval Parameters and Prompt Design}
%
%Profession retrieval was performed using the template:
%
%\[
%\texttt{"A photo of a/an \{profession\}"}
%\]
%
%This phrasing was selected to encourage naturalistic depictions of individuals rather than symbolic or abstract representations.
%
%For each profession:
%\begin{itemize}
%	\item The top 125{,}000 images were initially retrieved.
%	\item A cosine similarity threshold of 0.15 was applied.
%	\item From this candidate pool, the top 1{,}000 highest-scoring structurally valid images were selected.
%\end{itemize}
%
%The similarity threshold was empirically chosen to remove near-random matches while maintaining a sufficiently large candidate pool to enable robust filtering.
%
%Additionally, the following exclusion keywords were applied during retrieval filtering: \emph{Cartoon, NSFW, Sex, Naked, Object, Sign, Logo, Document, Paper, Page}. These exclusions were introduced after exploratory analysis revealed frequent retrieval of document-based or non-photographic imagery (e.g., legal contracts for ``lawyer''), which are unsuitable for demographic analysis.
%
%\subsection{Sample Size Justification}
%
%For each of the 200 professions, 1{,}000 images were retained, yielding a balanced subset of 200{,}000 images.
%
%Under a binomial assumption for spurious feature presence (feature present vs.\ absent), a sample size of 1{,}000 per profession yields a worst-case 95\% margin of error of approximately $\pm 3.1\%$ for proportion estimates (maximal at $p = 0.5$). This ensures that observed differences in feature prevalence across professions are unlikely to arise from sampling noise.
%
%Selecting an equal number of images per profession produces a balanced evaluation set. This design prevents more frequently occurring professions in the underlying dataset from disproportionately influencing aggregate statistics, thereby enabling controlled cross-profession comparison. When aggregated across all professions, the 200{,}000-image subset yields highly precise overall prevalence estimates (worst-case $\pm 0.22\%$ at 95\% confidence), though such aggregate measures reflect an equal-weighted average rather than the natural profession distribution within LAION-5B.
%
%\subsection{Human-Centric Structural Filtering}
%
%CLIP retrieval does not guarantee that returned images contain clearly identifiable individuals. Therefore, a multi-stage structural validation pipeline was applied using YOLO-based detection models and MediaPipe FaceMesh segmentation.
%
%Filtering was performed using:
%\begin{itemize}
%	\item \texttt{yolov8n-face.pt} (face detection),
%	\item \texttt{yolov8n.pt} (person detection).
%\end{itemize}
%
%Face detection required:
%\begin{itemize}
%	\item Exactly one detected face,
%	\item Confidence $\geq 0.6$.
%\end{itemize}
%
%Person detection required:
%\begin{itemize}
%	\item At least one detected person (class ID = 0),
%	\item Confidence $\geq 0.5$.
%\end{itemize}
%
%To ensure that detected faces corresponded to detected persons, spatial consistency constraints were enforced:
%\begin{itemize}
%	\item Face centroid must lie within a person bounding box \textbf{or}
%	\item Intersection-over-Union (IoU) $\geq 0.3$.
%\end{itemize}
%
%Additionally, the face-to-person area ratio was required to exceed 0.02 to prevent degenerate detections. Following successful validation, MediaPipe FaceMesh segmentation was applied to extract a refined facial region. Images failing any stage were marked invalid.
%
%Model selection was informed by empirical evaluation on the WIDER Face validation dataset (see Table~\ref{add_table}), where YOLOv8-face demonstrated favourable performance relative to MediaPipe Face Detection, SCRFD, and RetinaFace. Person detection was included as an additional structural safeguard to prevent false-positive face detections in non-human contexts.
%
%\subsection{Filtering Outcomes and Acceptance Rate Distribution}
%
%Across all professions, a total of 25{,}125{,}679 candidate images were processed following initial retrieval. Of these, 8{,}420{,}684 images satisfied all structural acceptance criteria, yielding an overall acceptance rate of 33.51\%.
%
%Acceptance rates varied substantially across professions. The mean acceptance rate was 33.51\%, the median 33.67\%, and the standard deviation 10.71\%. The lowest observed acceptance rate was 9.55\%, while the highest was 59.21\%.
%
%Figure~\ref{fig:acceptance_histogram} presents the distribution of acceptance rates across professions. The histogram demonstrates that structural validity following CLIP retrieval is highly profession-dependent. Certain professions systematically yield retrieval results that more frequently depict clear single-person imagery, while others produce noisier or object-dominated results. This variability indicates that structural filtering does not affect all professions uniformly and may introduce differential pruning effects across categories.
%
%The overall rejection rate of approximately 66\% further illustrates that similarity-based retrieval alone is insufficient for ensuring demographic suitability. Explicit structural validation is therefore necessary prior to fairness or spurious feature analysis.
%
%\subsection{Duplicate Handling}
%
%Images with identical embeddings were removed to prevent duplication within the profession-conditioned subset. Cross-profession duplication was theoretically possible if an image ranked highly for multiple profession prompts; however, such instances were empirically rare and were retained when present, as they represent semantically relevant examples for each respective query.
%
%\subsection{Random Control Subset}
%
%To evaluate whether identified spurious features were specific to profession-conditioned retrieval or reflective of broader dataset characteristics, an additional control subset of 10{,}000 randomly sampled images was constructed. These images were drawn without replacement from the processed ReLAION subset, excluding images used in the profession-conditioned dataset.
%
%Sampling was seed-controlled to ensure reproducibility. This control group serves to distinguish retrieval-induced bias from underlying dataset bias and to assess whether observed spurious correlations generalise beyond occupation-aligned subsets.
%
%In line with the goals of this reaserch paper generated images derived from models trained on this dataset were required. Several models were considered and tested these being SDv1.5, SDv1.5 + Speed Lora, SD v2.1, SDXL, SDXL+Speed Lora, SDXL-Turbo, SSD-1B,  SSD-1B+Speed Lora, PixArt-Alpha, PixArt-Alpha LCM, Segmind Vega + Vega RT Lora. These models were compared to one anotehr in terms of inference speed, image quality was not considered in relation to identifying the model to utilise as this is not relevant to the question tackled in this thesis, serving only as a secondary factor for model selection. In terms of inference time as per the table \ref{tab:generative-model-inference-time} the SDv1.5 + Speed Lora model was chosen as it was amongst the fastest models whilst producing coherent images.



%\subsection{Image Generation Pipeline}
%
%In order to analyse representational bias at the generative stage, synthetic images were produced using diffusion-based text-to-image models trained on LAION-scale datasets. Multiple models were evaluated, including SDv1.5, SDv1.5 + Speed LoRA, SDv1.2, SDXL, SDXL + Speed LoRA, SDXL-Turbo, SSD-1B, SSD-1B + Speed LoRA, PixArt-Alpha, PixArt-Alpha LCM, and Segmind Vega (with Vega RT LoRA). 
%
%Models were compared primarily on inference speed rather than perceptual image quality. Since the objective of this thesis is to analyse structural and demographic bias rather than aesthetic fidelity, image quality was considered a secondary factor. As shown in Table~\ref{tab:generative-model-inference-time}, SDv1.5 with Latent Consistency Model (LCM) LoRA was selected. This configuration achieved one of the lowest inference times while maintaining structurally coherent outputs suitable for downstream validation.
%
%Stable Diffusion v1.5 was loaded with LCM LoRA weights and fused for fast sampling using a four-step inference schedule. The scheduler was replaced with an LCM scheduler to enable low-step generation while maintaining semantic consistency. Generation was performed in half-precision (FP16) on GPU.
%
%\subsection{Prompt Design and Justification}
%
%The final prompt used for profession-conditioned generation was:
%
%\begin{quote}
%	\texttt{"full-body realistic photo of a \{profession\}, standing, centered, in a professional setting appropriate for their occupation, sharp focus, 35mm lens, natural professional lighting"}
%\end{quote}
%
%Each component was selected for functional and fairness-related reasons.
%
%\paragraph{Full-body framing.}
%Diffusion models trained on LAION-scale data exhibit a strong bias toward portrait-style close-ups. Explicitly requesting ``full-body'' reduces framing bias and ensures compatibility with structural validation constraints requiring a full detected person bounding box.
%
%\paragraph{Standing, centered.}
%Explicit pose and spatial positioning tokens reduce multi-subject ambiguity and mitigate edge-cropping artefacts. This increases the probability of successful face–person spatial association during validation.
%
%\paragraph{Professional setting appropriate for their occupation.}
%Rather than specifying explicit environments (e.g., hospital, courtroom), which are known to amplify demographic priors, an abstract contextual phrase was used. This preserves semantic grounding while limiting overfitting to stereotypical scenes.
%
%\paragraph{Sharp focus.}
%Low-step LCM sampling can introduce blur. High-frequency detail cues such as ``sharp focus'' improve edge clarity and increase downstream detector reliability.
%
%\paragraph{35mm lens.}
%Camera tokens act as implicit geometric priors influencing subject distance and body proportions. A 35mm focal length promotes natural perspective and full-body framing without excessive distortion.
%
%\paragraph{Natural professional lighting.}
%Lighting significantly influences aesthetic regression scores and object detection reliability. Neutral professional lighting reduces noise artefacts and improves face and person detection confidence.
%
%\subsection{Generation with Structural Validation}
%
%Images were generated for each of the 200 professions using the above prompt template. For each profession, 1,000 valid images were required.
%
%The generation pipeline incorporated asynchronous structural validation (see \texttt{ImageGeneration.py}~\cite{turn2file0}):
%
%\begin{itemize}
%	\item YOLO person detection (confidence $\geq 0.4$),
%	\item YOLO face detection (confidence $\geq 0.4$, exactly one face required),
%	\item Spatial face–person consistency via centroid inclusion or bounding box expansion,
%	\item Face-to-person area ratio thresholding,
%	\item MediaPipe FaceMesh segmentation for anatomical refinement,
%	\item CLIP-based aesthetic scoring (minimum score threshold = 3.0).
%\end{itemize}
%
%An image was accepted only if all criteria were satisfied. Images failing any stage were stored in an invalid directory with the rejection reason encoded in the filename.
%
%\subsection{Generation Without Validation (Control)}
%
%To assess generative behaviour absent structural filtering, an additional control set of 10,000 images was generated using blank prompts and no validation constraints (see \texttt{ImageGeneration\_NoValidation.py}~\cite{turn2file1}). 
%
%These images were sampled using identical inference parameters but without YOLO-based filtering, spatial constraints, or aesthetic thresholding. This control group enables comparison between structurally constrained and unconstrained generative outputs, isolating the effect of validation logic on observed representational patterns.
%
%\subsection{Rationale for Structural Validation in Generation}
%
%Unlike retrieval from ReLAION-5B, generative models can produce:
%\begin{itemize}
%	\item Multiple individuals,
%	\item Truncated bodies,
%	\item Disconnected faces,
%	\item Anatomically inconsistent compositions,
%	\item Low-quality or blurred renderings.
%\end{itemize}
%
%Since downstream demographic and spurious feature analysis requires clearly identifiable single individuals, structural validation is necessary to ensure comparability between retrieved and generated subsets. 
%
%This design ensures that differences observed between real and synthetic images reflect model behaviour rather than artefacts of composition, cropping, or detection instability.
%
%\subsection{Structural Acceptance Rates for Generated Images}
%
%Following generation and validation, acceptance rates were computed per profession in an identical manner to the retrieval analysis. 
%
%Across 200 professions, the mean structural acceptance rate for generated images was 66.38\%, with a median of 71.00\% and standard deviation of 16.60\%. Acceptance rates ranged from approximately 10\% to over 90\%.
%
%Figure~\ref{fig:sd_acceptance_histogram} illustrates the distribution of acceptance rates across professions. Compared to retrieval from ReLAION-5B (mean acceptance 33.51\%), generated images exhibit substantially higher structural validity. This indicates that diffusion outputs, when guided by explicit full-body and framing constraints, are more likely to satisfy single-person detection and spatial consistency requirements than real-world web imagery.
%
%However, variability across professions remains pronounced. Certain occupations consistently produce high acceptance rates, suggesting that the diffusion model readily associates these roles with canonical single-person depictions. In contrast, other professions exhibit lower structural acceptance, potentially reflecting greater compositional ambiguity or weaker learned visual priors within the model.
%
%The higher mean acceptance rate suggests that structural constraints imposed via prompt engineering and diffusion priors systematically bias generation toward detector-compatible compositions. This difference between retrieved and generated subsets must be considered when comparing downstream demographic or spurious feature prevalence.



\graphicspath{{content/chapters/3_methodology/figures/}}

\chapter{Methodology}%
\label{chp:methodology}


\section{Image Retrieval}

\subsection{Data Acquisition and Subsetting}

Re-LAION-5B is distributed as 128 shards, each approximately 3.62\,GB in size, together comprising the full corpus of URL--caption pairs. Due to storage and computational constraints, processing was restricted to four shards, yielding an initial subset of 65,552,809 URL--caption pairs. Shards were selected uniformly at random from the full index. Because LAION shards are generated during distributed web crawling without semantic ordering, shard index does not correspond to content category, domain, or temporal structure; random selection therefore provides a principled basis for treating the subset as representative of the broader corpus. Under a binomial sampling model, a worst-case 95\% margin of error of approximately $\pm 0.04\%$ is achieved for proportion estimates at this scale, confirming that sampling variability alone is unlikely to account for observed distributional patterns. Nonetheless, the presence of unknown crawl-time biases---such as domain saturation effects arising from crawl order---cannot be excluded without access to LAION crawl metadata. All findings are therefore reported at the computed confidence level and should be interpreted as representative estimates rather than exact population parameters.

\subsection{Image Source Distribution}

Analysis of valid image origins revealed that 18.79\% of successfully retrieved images originated from five domains: \texttt{cdn.shopify.com} (3,053,782 images), \texttt{i.pinimg.com} (2,585,519), \texttt{i.ebayimg.com} (966,820), \texttt{images-na.ssl-images-amazon.com} (886,054), and \texttt{thumbs.dreamstime.com} (678,516). This concentration indicates strong domain centralisation within the subset, with large commercial platforms contributing disproportionately to the retrieved image pool. As a consequence, product imagery and stock photography may be overrepresented relative to content from smaller independent domains, and this skew should be considered when interpreting bias estimates derived from retrieved imagery.

\subsection{Image Validity Criteria}

Since Re-LAION-5B distributes URL--caption pairs rather than image binaries, each URL was resolved and the corresponding image downloaded. An image was deemed valid only if it satisfied all of the following conditions.

\paragraph{HTTP Resolution.} A GET request must complete successfully with an HTTP 2xx status code within a 10-second timeout. Requests were issued with a browser-like \texttt{User-Agent} header and a domain-derived \texttt{Referer} header to reduce server-side rejection. Any network exception---including timeouts, DNS failures, SSL errors, and connection errors---resulted in immediate discard.

\paragraph{Image Decoding.} The response body must be successfully opened via \texttt{PIL.Image.open()} and fully loaded via \texttt{image.load()}, ensuring the file is neither corrupted nor truncated.

\paragraph{RGB Conversion.} Images must be successfully converted to RGB. Non-RGB formats (e.g., RGBA, grayscale) were normalised to RGB; images that failed conversion were discarded.

\paragraph{Batch-Level Embedding Validation.} If embedding computation failed for a batch, recursive binary splitting was applied: the batch was subdivided into two equal halves and each half re-attempted independently. This process continued recursively until either a valid sub-batch was processed or a batch of size one failed, at which point that individual image was excluded. This approach avoids falling back to fully serial processing and substantially reduces recovery overhead for large batches containing isolated corrupted images.

\subsection{Decay Rate}

From 65,552,809 attempted URLs, 43,491,938 images were successfully retrieved and validated, yielding a decay rate of:

\[
\text{Decay Rate} = 1 - \frac{43{,}491{,}938}{65{,}552{,}809} \approx 0.337
\]

This level of link attrition is consistent with prior observations of web-scale dataset degradation and reflects natural URL expiry over time. However, link decay is not necessarily uniform: if certain domains or content types become unavailable at disproportionate rates, the surviving image pool may be systematically skewed relative to the original dataset composition, introducing survivorship bias that cannot be quantified post-hoc.

\subsection{Image Preprocessing}

All images were preprocessed using the canonical pipeline associated with \texttt{openai/clip-vit-large-patch14}, accessed via the HuggingFace \texttt{AutoProcessor}. Preprocessing was applied deterministically and included resizing to 224$\times$224 pixels with aspect ratio preservation prior to centre cropping, followed by pixel rescaling and channel-wise normalisation using CLIP-specific mean and standard deviation values. No data augmentation was applied. Output tensors were of type \texttt{float32}.

\subsection{Embedding Computation}

CLIP image embeddings were recomputed for all valid images, as the original LAION-5B embeddings are no longer publicly available. Embeddings were extracted using \texttt{openai/clip-vit-large-patch14} (embedding dimensionality: 768; patch resolution: 14$\times$14), pretrained on OpenAI's WebImageText corpus.

The download and embedding pipeline employed a producer--consumer architecture to maximise throughput. A dedicated download thread retrieved URL batches in parallel using a \texttt{ThreadPoolExecutor} with 100 concurrent workers, issuing batches of 200 images to an in-memory queue with a capacity of 10 batches to prevent unbounded memory accumulation. A consumer thread dequeued each completed download batch and forwarded valid images to the CLIP encoder for embedding extraction via \texttt{model.get\_image\_features()} under \texttt{torch.no\_grad()} in evaluation mode, with GPU acceleration where available.

Each embedding vector $\mathbf{e}$ was L2-normalised:

\[
\mathbf{e}_{\mathrm{norm}} = \frac{\mathbf{e}}{\|\mathbf{e}\|_2}
\]

This enforces unit-length representations and enables cosine similarity retrieval via inner-product equivalence.

Results were accumulated in an in-memory buffer and flushed to Parquet files incrementally. A checkpoint was written every 5,000 processed entries; at most three checkpoints were retained at any time, with older checkpoints deleted automatically to manage disk usage. Upon reaching 2,000,000 entries, a checkpoint was promoted to a finalised part file and the buffer reset. This design ensured fault tolerance and resumability across the multi-month processing window without requiring full reprocessing after interruption. The pipeline additionally incorporated internet connectivity monitoring with exponential backoff, pausing download activity during transient network outages and resuming automatically upon restoration.

\subsection{Model Selection Justification}

While the original LAION-5B filtering pipeline employed CLIP ViT-B/32 embeddings, this work uses the higher-capacity ViT-L/14 variant. ViT-L/14 offers 768-dimensional embeddings, finer spatial granularity (14$\times$14 patches versus 32$\times$32), and stronger semantic alignment in retrieval benchmarks. Since the objective of this work is bias analysis rather than replication of LAION filtering thresholds, the use of a higher-capacity backbone is methodologically appropriate. It should be noted that all embeddings inherit inductive biases from their pretraining corpus, and retrieval behaviour remains model-dependent.

\subsection{Computational Environment}

The download and embedding pipeline was executed on an NVIDIA RTX 3060 Ti GPU with 64\,GB of system RAM. Processing spanned approximately two to three months. The extended duration may have introduced minor temporal drift in URL availability, as pages accessible at the start of processing may have become unavailable by its conclusion.

\subsection{Ethical Considerations}

\begin{itemize}
	\item Images were accessed exclusively via publicly available URLs.
	\item No authentication or access restrictions were bypassed.
	\item Images were used strictly for academic research purposes.
	\item No redistribution of image binaries was performed.
	\item Analysis was conducted at aggregate statistical levels.
	\item No attempt was made to identify individuals.
	\item No facial recognition or re-identification techniques were applied.
	\item Where possible, only embeddings were retained for analytical tasks.
	\item The research objective is to analyse bias rather than amplify or deploy it.
\end{itemize}

\subsection{Limitations}

\begin{itemize}
	\item Only 4 of 128 shards were processed; findings apply strictly to the processed subset.
	\item Unknown crawl-time biases cannot be excluded without access to LAION crawl metadata.
	\item Approximately 33.7\% of URLs were unavailable, introducing potential survivorship bias.
	\item 18.79\% of valid images originated from five domains, indicating domain concentration.
	\item Embeddings reflect biases inherent to the CLIP pretraining corpus.
	\item Image resizing to 224$\times$224 may reduce fine-grained demographic cues.
	\item A 10-second request timeout may classify slow-responding servers as invalid.
	\item No geographic replication of downloads was performed.
	\item Results are encoder-dependent and not model-agnostic.
	\item No additional content filtering was applied prior to embedding extraction.
\end{itemize}
















\graphicspath{{content/chapters/3_methodology/figures/}}

%% ============================================================
%% Section: Profession-Conditioned Subset Construction
%% ============================================================

\section{Profession-Conditioned Image Retrieval}

\subsection{Profession-Conditioned Subset Construction}

Following large-scale image retrieval and embedding computation, a structured subset of images was required upon which spurious feature analysis would be conducted. In line with prior work examining occupational representation and stereotype formation in web-scale datasets (Section~\ref{subsec:occupations}), the analytical focus was placed on profession-conditioned imagery. A set of 200 diverse professions was constructed and aligned with ISCO-08 occupational categories to ensure systematic coverage across skill levels and domain areas. For completeness, the full profession list is provided in the Appendix.

\subsection{FAISS Index Construction}

To enable efficient similarity-based retrieval over the full embedding corpus of 43 million L2-normalised CLIP ViT-L/14 image embeddings, a dedicated FAISS index was constructed for each processed shard. Separate \texttt{IndexFlatIP} indexes were built per shard, with each index accompanied by a position-to-path mapping stored as a serialised Python list. Since embeddings were L2-normalised during extraction, inner-product search over \texttt{IndexFlatIP} is equivalent to cosine similarity retrieval.

Index construction processed each Parquet file row-group by row-group, parsing embeddings in their stored fixed-size list format and adding them to the index in batch. Both the FAISS index and the path mapping were checkpointed periodically using atomic file replacement to ensure crash recovery without data corruption. A JSON state file tracked the number of row groups processed per Parquet file, enabling full resumability after interruption. At inference time, shard indexes were streamed sequentially, with per-shard retrieval results aggregated and globally re-ranked by similarity score.

\subsection{Prompt Design}

The retrieval template used for each profession was:

\[
\texttt{"A photo of a/an \{profession\}"}
\]

The grammatical article was assigned programmatically: \textit{an} was selected when the profession name began with a vowel; \textit{a} otherwise. This rule is sufficient to produce grammatically correct CLIP prompts for all occupation names in the profession list. Gender-neutral phrasing was used throughout to avoid reinforcing gender priors at the query stage.

\subsection{Negative Prompt Dominance Filtering}

To suppress retrieval of semantically irrelevant content, a negative prompt dominance filter was applied during search. For each candidate image, its CLIP embedding was compared against a set of negative concept embeddings derived from the following prompts: \emph{Cartoon, NSFW, Sex, Naked, Clothing, Object, Sign, Logo, Document, Paper, Page, Animal}. These exclusions were motivated by exploratory analysis, which revealed frequent retrieval of document imagery (e.g., legal contracts for \emph{Lawyer}), symbolic or non-photographic content (e.g., logos and signs), and non-human subjects.

A candidate image was rejected if the maximum cosine similarity between its embedding and any negative concept embedding satisfied:

\[
\text{sim}_{\text{neg}} \geq \text{sim}_{\text{pos}} - \delta
\]

where $\text{sim}_{\text{pos}}$ denotes the cosine similarity to the target profession prompt and $\delta = 0.01$ is a dominance margin. This condition discards images where a negative concept is either more semantically similar than the target or within a margin of 0.01 of the target score, while retaining images where the positive prompt is clearly dominant.

\subsection{Embedding-Level Deduplication}

To prevent the same image from appearing multiple times within the result set for a given profession, embedding-level deduplication was applied during retrieval. Each candidate image embedding was reconstructed from the FAISS index and hashed via its byte representation. Embeddings whose hash had already been observed for the current query were discarded prior to result accumulation. This deduplication is performed within a single profession query; cross-profession deduplication was not enforced, as the same image may legitimately be semantically relevant to multiple professions.

\subsection{Retrieval Parameters}

For each of the 200 professions, the top 125{,}000 candidates were retrieved per shard and aggregated globally. A cosine similarity threshold of 0.15 was applied to remove near-random matches while maintaining a sufficiently large candidate pool for downstream structural filtering. Following YOLO-based structural validation (Section~\ref{subsec:yolo_filtering_retrieval}), the top 1{,}000 highest-scoring structurally valid images were retained per profession. This two-stage design decouples semantic relevance (FAISS retrieval) from structural suitability (YOLO filtering), allowing the structural stage to prune without pre-restricting the similarity pool.

\subsection{Human-Centric Structural Filtering}
\label{subsec:yolo_filtering_retrieval}

CLIP-based retrieval does not guarantee that retrieved images contain clearly identifiable individuals. A multi-stage structural validation pipeline was therefore applied to enforce human-centric image suitability.

Images were loaded in batches of 128 using a parallel loading stage that prioritised TurboJPEG for JPEG-format inputs, with OpenCV as a fallback. Batch-level inference was performed with YOLO models operating in half-precision (FP16) to maximise throughput. The two detection models used were:

\begin{itemize}
	\item \texttt{yolov8n-face.pt} for face detection,
	\item \texttt{yolov8n.pt} for whole-body person detection.
\end{itemize}

Face detection required exactly one detected face with confidence $\geq 0.6$. Person detection required at least one detected person (YOLO class ID = 0) with confidence $\geq 0.5$. To verify that the detected face corresponded to one of the detected persons, a spatial consistency check was applied: the face centroid must lie within a person bounding box, or the intersection-over-union (IoU) between the face and person bounding boxes must be $\geq 0.3$. Additionally, the ratio of face area to the associated person bounding box area was required to exceed $0.02$, preventing degenerate detections in which a small false-positive face detection is spuriously associated with a large person region.

Images satisfying person and face detection criteria but for which no person bounding box was spatially consistent with the detected face were still admitted, provided that MediaPipe FaceMesh segmentation succeeded on the face region. This fallback ensures that images in which a face is clearly detectable but the associated person bounding box is slightly misaligned are not needlessly discarded.

Following successful detection, MediaPipe FaceMesh (configured with \texttt{static\_image\_mode=True}, \texttt{max\_num\_faces=1}, and \texttt{refine\_landmarks=True}) was applied to the face region of interest. FaceMesh landmarks were used to compute a convex hull mask, from which a tightly cropped face ROI was extracted with a 10-pixel padding. Images for which FaceMesh failed to return landmarks were discarded. Accepted images and their corresponding face crops were written asynchronously using a parallel thread pool executor.

Progress was tracked using a completion manifest: per-profession completion status was persisted atomically to a JSON manifest file after each profession group was fully processed, enabling resumability without reprocessing completed groups.

\subsection{Filtering Outcomes and Acceptance Rate Distribution}

Across all 200 professions, 25{,}125{,}679 candidate images were submitted to structural validation following initial FAISS retrieval. Of these, 8{,}420{,}684 images satisfied all criteria, yielding an overall acceptance rate of 33.51\%.

Acceptance rates varied substantially across professions. The mean acceptance rate was 33.51\%, the median 33.67\%, and the standard deviation 10.71\%. The minimum observed acceptance rate was 9.55\% and the maximum 59.21\%.

Figure~\ref{fig:acceptance_histogram} illustrates the distribution of per-profession acceptance rates across the 200 evaluated professions. The distribution is approximately unimodal and centred around the mean, though with a visible left tail indicating that a subset of professions consistently produce structurally unsuitable retrieval results. This variability is consistent with the expectation that CLIP retrieval quality is profession-dependent: occupations associated with clear single-person visual conventions (e.g., \emph{Chef}, \emph{Nurse}) are likely to yield more structurally valid results than those whose online visual representation is dominated by equipment, environments, or group scenes.

The overall rejection rate of approximately 66\% demonstrates that CLIP-based retrieval alone is insufficient for demographic analysis. Explicit structural validation is therefore a necessary preprocessing stage prior to spurious feature or fairness analysis.

%\begin{figure}[htbp]
%	\centering
%	\includegraphics[width=0.9\textwidth]{ReLaion-Distribution.png}
%	\caption{Distribution of image acceptance rates across professions following CLIP retrieval and YOLO-based structural filtering of Re-LAION-5B images. Mean: 33.51\%, Median: 33.67\%, Std Dev: 10.71\%.}
%	\label{fig:acceptance_histogram}
%\end{figure}

\subsection{Random Control Subset}
\label{subsec:random_control}

A random control subset of 10{,}000 images was constructed from the four processed Re-LAION-5B shards, with explicit exclusion of any image already assigned to the profession-conditioned dataset. Unlike the profession-conditioned subset, which was retrieved using occupation-specific CLIP queries and subsequently filtered to retain only images depicting a single clearly identifiable person, the control subset was drawn without any semantic or structural constraint. Images in this subset may depict people, objects, scenes, or any other content present in the underlying corpus.

The primary purpose of this subset is to assess whether spurious low-level features identified in the profession-conditioned imagery generalise beyond images depicting people. If the same features are present at comparable rates in images drawn without any human-centric retrieval or filtering constraint, this strengthens the argument that those features are properties of the Re-LAION-5B dataset as a whole rather than artefacts introduced by the CLIP retrieval process, the negative dominance filtering, or the YOLO-based structural validation pipeline.

Exclusion of profession-conditioned images was enforced using a bitset-based strategy. A streaming pass over the profession-conditioned retrieval JSONL file populated per-shard bitsets indexed by integer image filename (e.g., \texttt{1234567.jpg}), marking each used image index as unavailable. The bitset representation required only a few megabytes of memory per shard ($\mathcal{O}(\texttt{max\_index} / 8)$ bytes), making the approach scalable to shard sizes of up to 17 million images without loading the full index into memory. Random sampling then drew uniformly from the remaining available indices. Sampling was performed with a fixed random seed (seed = 42) to ensure full reproducibility.

\subsection{Sample Size Justification}

For each of the 200 professions, 1{,}000 images were retained following structural validation, yielding a balanced profession-conditioned subset of 200{,}000 images. Under a binomial model for binary feature presence, a per-profession sample of $n = 1{,}000$ yields a worst-case 95\% margin of error of:

\[
\text{MOE} = z_{0.975} \cdot \sqrt{\frac{p(1-p)}{n}} = 1.96 \cdot \sqrt{\frac{0.25}{1000}} \approx \pm 3.1\%
\]

This margin is sufficiently small to detect meaningful cross-profession differences in spurious feature prevalence. When aggregated across all 200 professions, the full 200{,}000-image subset yields highly precise overall estimates (worst-case $\pm 0.22\%$ at 95\% confidence). Equal sampling across professions produces a balanced evaluation set that prevents high-frequency professions from disproportionately influencing aggregate statistics, enabling controlled cross-profession comparison. This design does, however, imply that aggregate prevalence estimates reflect an equal-weighted average over professions rather than the natural profession frequency distribution within Re-LAION-5B.








\graphicspath{{content/chapters/3_methodology/figures/}}

%% ============================================================
%% Section: Image Generation Pipeline
%% ============================================================

\section{Image Generation Pipeline}

\subsection{Model Selection}

To analyse representational bias at the generative stage, profession-conditioned synthetic images were produced using diffusion-based text-to-image models trained on LAION-scale datasets. Eleven model configurations were evaluated: Stable Diffusion v1.5 (SDv1.5), SDv1.5 with LCM LoRA, Stable Diffusion v1.2, SDXL Base, SDXL with LCM LoRA, SDXL-Turbo, SSD-1B, SSD-1B with LCM LoRA, PixArt-Alpha, PixArt-Alpha LCM, and Segmind Vega with VegaRT LoRA. All models were evaluated at 512$\times$512 pixel resolution using a fixed profession prompt for a Doctor.

Model comparison was based primarily on inference throughput rather than perceptual image quality. Since the objective of this thesis is to analyse structural and demographic bias rather than aesthetic fidelity, image quality was considered only as a secondary factor sufficient to ensure structurally coherent outputs suitable for downstream detection-based validation. As shown in Table~\ref{tab:generative-model-inference-time}, SDv1.5 with LCM LoRA achieved one of the lowest per-image inference times while producing images with sufficient structural coherence.

SDv1.5 was loaded with the \texttt{latent-consistency/lcm-lora-sdv1-5} adapter and fused for fast inference. LoRA weights were fused into the base model before generation, eliminating adapter overhead at inference time. The default DDPM scheduler was replaced with \texttt{LCMScheduler}, which enables high-quality generation at substantially reduced step counts through consistency distillation. A four-step inference schedule was used for all validated generation runs; the testing notebook confirmed that fewer than four steps introduces visual distortion, while more than four yields no measurable improvement in structural quality. Generation was performed in half-precision (FP16) on GPU, with TF32 enabled for matrix multiplications and cuDNN operations to maximise throughput without sacrificing numerical precision.

\subsection{Prompt Design and Justification}

The prompt template used for all profession-conditioned generation was:

\begin{quote}
	\texttt{full-body realistic photo of a \{profession\}, standing, centered, in a professional setting appropriate for their occupation, sharp focus, 35mm lens, natural professional lighting}
\end{quote}

Each component was chosen for a specific functional or fairness-related reason.

\paragraph{Full-body framing.}
Stable Diffusion models trained on LAION-scale data exhibit a strong prior toward portrait-style close-ups. Explicitly specifying full-body framing overrides this default and ensures compatibility with downstream structural validation, which requires a full detected person bounding box to confirm subject completeness.

\paragraph{Standing, centered.}
Explicit pose and spatial positioning tokens reduce compositional ambiguity arising from multiple subjects or edge-cropping artefacts. These constraints increase the probability that the detected face falls within the upper half of the person bounding box, satisfying the anatomical plausibility check applied during structural validation.

\paragraph{Professional setting appropriate for their occupation.}
Rather than specifying concrete environments (e.g., hospital, courtroom), which are known to amplify demographic stereotypes by conditioning the model on environment-associated appearance priors, an abstract contextual phrase was used. This preserves profession-relevant semantic grounding while limiting activation of environment-specific demographic shortcuts.

\paragraph{Sharp focus.}
The four-step LCM sampling schedule can introduce diffusion-specific blur. High-frequency cues such as sharp focus improve edge fidelity, increasing the reliability of downstream face and person detection.

\paragraph{35mm lens.}
Camera focal-length tokens act as implicit geometric priors, influencing subject distance and body proportions. A 35mm specification promotes naturalistic perspective and full-body framing without barrel or perspective distortion.

\paragraph{Natural professional lighting.}
Lighting conditions significantly influence CLIP-based aesthetic scores and detection confidence. Neutral professional lighting tokens suppress noise artefacts and promote image quality scores that meet the downstream aesthetic filtering threshold.

\subsection{Guidance Scale and Its Methodological Implications}

The classifier-free guidance scale was fixed at $s = 1.0$ across all validated generation runs. At $s = 1$, no guidance amplification is applied, and the model produces outputs conditioned purely on its learned prior without any prompt-signal exaggeration.

Standard Stable Diffusion deployment uses guidance scales in the range $s = 7$--$7.5$. This choice was deliberate. Classifier-free guidance is an inference-time amplification mechanism that scales the difference between conditional and unconditional predictions; elevated guidance increases prompt adherence but also amplifies any demographic or compositional bias encoded in the learned prior. By fixing $s = 1$, the generated distribution reflects the model's intrinsic learned associations rather than inference-time amplifications of those associations. As a consequence, any demographic bias identified under $s = 1$ represents a conservative lower bound: at guidance scales used in production deployment, the same biases would be further amplified. If bias is detectable at its minimum---when guidance amplification is absent---it follows that the bias is structural and attributable to model weights and training data, not to inference-time configuration.

This stands in contrast to the unconditional control subset, which was generated with $s = 0.0s$ and an empty prompt, producing samples from the model's learned prior without any profession conditioning. The guidance scale argument is therefore specific to the profession-conditioned runs; the control subset serves a different purpose and is not subject to the same inference-time amplification concern.

\subsection{Asynchronous Generation and Validation Architecture}

Images were generated for each of the 200 professions in batches of 2. For each profession, generation continued until 1{,}000 structurally valid images had been accepted. To avoid blocking GPU generation during CPU-bound validation, structural validation was decoupled from generation using an asynchronous producer--consumer architecture.

A \texttt{ValidationWorker} pool was instantiated with a configurable number of worker threads (default: 2). Each worker thread independently loaded its own copies of the YOLO person detection model, YOLO face detection model, and MediaPipe FaceMesh. Validation tasks were submitted to a shared task queue and results returned via a result queue, allowing the main GPU generation loop to remain active while CPU workers processed previously generated images in parallel. Generated images were held in an in-memory cache keyed by unique image identifier, and retrieved for saving or discarding once the corresponding validation result arrived. File I/O was performed asynchronously via a separate thread pool executor with 4 workers to avoid blocking either generation or validation.

\subsection{Structural Validation Criteria}

An image was accepted only if it satisfied all of the following conditions.

\paragraph{Person detection.} The YOLO person detection model (\texttt{yolo12s.pt}) must detect at least one bounding box with class ID 0 and confidence $\geq 0.4$.

\paragraph{Face detection.} The YOLO face detection model (\texttt{yolov12l-face.pt}) must detect exactly one face with confidence $\geq 0.4$. Images with zero or multiple detected faces were rejected.

\paragraph{Face crop validity.} The face bounding box must define a non-empty region of interest. Degenerate crops of zero area were rejected.

\paragraph{FaceMesh segmentation.} MediaPipe FaceMesh (configured with \texttt{static\_image\_mode=True}, \texttt{max\_num\_faces=1}, and \texttt{refine\_landmarks=True}) must successfully return landmarks for the face region of interest, with a minimum of 10 landmark points. Landmarks were used to compute a convex hull mask over the face, from which a tightly cropped segmented face ROI was extracted with a 10-pixel boundary padding.

\paragraph{Spatial and anatomical consistency.} For each detected person bounding box, a composite spatial score was computed:

\[
\text{score} = 10 \cdot \mathbf{1}[\text{centroid\_inside}] + 2 \cdot \mathbf{1}[\text{upper\_half}] + \min\!\left(\frac{A_\text{face}}{A_\text{person}},\, 0.05\right) \cdot 20
\]

where $\mathbf{1}[\text{centroid\_inside}]$ indicates that the face centroid lies within the person bounding box, $\mathbf{1}[\text{upper\_half}]$ indicates that the face centroid lies in the upper half of the person bounding box (an anatomical plausibility constraint), and $A_\text{face}/A_\text{person}$ is the face-to-person area ratio. The person bounding box with the highest score was selected as the candidate association. If the face centroid did not fall within the candidate bounding box, an expanded bounding box (expanded 5\% horizontally, 15\% vertically from the top edge) was tested. Images for which no valid person--face spatial association could be established were rejected. Additionally, the face-to-person area ratio must exceed $0.002$ to prevent association with degenerate detections.

\paragraph{Aesthetic quality.} Following successful structural validation, each accepted image was scored using a CLIP-based aesthetic regressor. Image features were extracted using \texttt{open\_clip} with the ViT-L-14 backbone pretrained on LAION-2B (\texttt{laion2b\_s32b\_b82k}), and projected to a scalar aesthetic score via a lightweight linear head (\texttt{sa\_0\_4\_vit\_l\_14\_linear.pth}) trained to predict human aesthetic preference ratings. Images with an aesthetic score below 3.0 were rejected.

Images passing all criteria were saved to disk along with their corresponding FaceMesh face crops. Rejected images were stored in a separate invalid directory with the rejection reason encoded in the filename, enabling post-hoc analysis of failure mode distributions. A CSV metadata file was written incrementally, recording image filename, profession label, aesthetic score, face bounding box, and person bounding box for each accepted image.

\subsection{Resumability}

The generation pipeline supported per-profession resumability. At the start of each profession, the pipeline counted existing valid images in the output directory. If the target count of 1{,}000 had already been reached, the profession was skipped without re-running generation. This design tolerated interruptions across the full 200-profession generation run without requiring any previously accepted images to be regenerated.

\subsection{Generation Without Validation: Unfiltered Control Subset}

A separate control set of 10{,}000 images was generated using blank prompts and no structural validation constraints. The inference configuration was otherwise identical to the validated profession-conditioned runs: Stable Diffusion v1.5 with LCM LoRA, four inference steps, and FP16 precision. The guidance scale was set to $s = 0.0$, yielding unconditional samples drawn from the model's learned prior with no prompt conditioning applied. As a consequence, these images carry no profession-specific semantic signal and are subject to none of the compositional constraints imposed by the generation prompt. They may depict people, objects, scenes, or abstract content, mirroring the unconstrained character of the Re-LAION-5B random control subset.

No YOLO-based person or face detection, spatial consistency checks, aesthetic thresholding, or FaceMesh segmentation were applied. Images were saved directly to disk as produced, preserving the raw generative distribution without any human-centric filtering.

The primary purpose of this subset is the same as that of the Re-LAION-5B random control: to assess whether spurious low-level features identified in the validated profession-conditioned images generalise to images generated without any human-centric constraint. If the same features appear at comparable rates in unconstrained generative outputs, this supports the conclusion that those features reflect properties of the model's learned prior rather than being introduced by the profession prompts or the structural validation pipeline. The parallel between the two control subsets --- one drawn from real web data, one sampled from a generative model --- also enables a direct comparison of the degree to which spurious correlations present in the training data are reproduced in model outputs.

\subsection{Rationale for Structural Validation in Generation}

Unlike retrieval from Re-LAION-5B, generative models produce images without guarantees of compositional suitability. Common failure modes include multiple subjects, truncated bodies, disconnected or misaligned faces, anatomically inconsistent compositions, and low-frequency blurred renderings arising from low-step inference. Since downstream demographic and spurious feature analysis requires clearly identifiable single individuals in consistent compositional configurations, structural validation is a necessary preprocessing stage. The parallel validation architecture ensures that this filtering does not constitute a bottleneck relative to GPU inference, as CPU-based detection and segmentation proceed concurrently with image generation.

The key design implication is that differences observed between real (Re-LAION-5B) and synthetic (Stable Diffusion) images in downstream analyses should reflect differences in model-learned content distributions rather than artefacts of composition, framing, or detection instability introduced by unfiltered generative output.

\subsection{Structural Acceptance Rates for Generated Images}

Following generation and validation, acceptance rates were computed per profession using the same methodology as the Re-LAION-5B retrieval analysis. Across all 200 professions, the mean structural acceptance rate for generated images was 66.38\%, with a median of 71.00\% and a standard deviation of 16.60\%. Acceptance rates ranged from approximately 10\% to over 90\%.

Figure~\ref{fig:sd_acceptance_histogram} illustrates the distribution of acceptance rates across professions. Compared to retrieval from Re-LAION-5B (mean acceptance 33.51\%), generated images exhibit substantially higher structural validity. This is consistent with the effect of the full-body prompt engineering, which biases the diffusion model toward compositions that satisfy single-person detection and face--person spatial association requirements. Nonetheless, per-profession variability remains substantial, indicating that certain occupations activate compositional priors within the model that systematically resist full-body single-subject framing. This variability may itself constitute a form of bias---professions for which the model more readily produces detector-compatible single-subject images may be more uniformly represented in downstream analyses than those with lower structural acceptance.

The systematic difference in mean acceptance rate between retrieved (33.51\%) and generated (66.38\%) subsets must be considered when comparing demographic or spurious feature prevalence across the two data sources, as the structural filtering stage does not operate identically across the two pipelines.


\section{}

%Methodology Note — CFG Scale Justification
%What to mention:
%State that CFG was fixed at s = 1 across all image generation experiments and explain what this means technically — that at s = 1 no guidance amplification is applied, so the model produces purely conditional outputs without any exaggeration of the prompt signal.
%Why to mention it:
%Reviewers familiar with SD v1.5 will notice s = 1 is atypical (the standard default is 7–7.5) and may flag it as a limitation if left unexplained. Getting ahead of this turns a potential weakness into a strength.
%How to frame it:
%Present it as a deliberate and principled methodological choice rather than a default. The reasoning is that fixing s = 1 isolates bias attributable to the model's learned training distribution from bias amplified at inference time. Since your thesis investigates training-data-derived spurious features, this is the most appropriate setting.
%The key argument to make:
%Any demographic bias identified under s = 1 represents a conservative lower bound. At the guidance scales used in real-world deployment (s = 7–7.5), the same biases would be amplified further. This strengthens rather than weakens your findings — if bias is detectable at its minimum, it will be more pronounced in practice, reinforcing the urgency of addressing it at the level of model weights and training data rather than inference settings.
%One line summary of the argument:
%Finding bias at s = 1 means the bias is structural, not incidental.

\section{Demographic Annotation}
\label{sec:demographic_annotation}

To enable the analysis of demographic bias in the Re-LAION-5B dataset and the images generated by Stable Diffusion v1.5, each face detected in the retrieved images required a gender label and a skin tone label. This section describes the evaluation pipeline used to select the best-performing model for each task, and the final annotation procedure applied to the full image sets.

% -----------------------------------------------------------
\subsection{Pre-processing}
\label{sec:annotation_preprocessing}

Before any model evaluation or final annotation was performed, a consistent face-segmentation step was applied to every image. Faces were isolated using MediaPipe FaceMesh~\cite{kartynnik2019realtime}, which fits a dense mesh to the detected face region and produces a precise segmentation mask. Pixels falling outside the mask were set to a uniform black background, yielding a cropped, background-removed face image that served as the input for all subsequent models.

This segmentation step was applied uniformly to both annotation tasks for two reasons. First, prior work on skin tone estimation has shown that models trained or evaluated on full-image crops are prone to exploiting background cues rather than true skin colour~\cite{wu2023skintone}; isolating the face region mitigates this confound. Second, applying the same pre-processing to both the gender classifier and the skin tone models ensures that both tasks operate on identical inputs, simplifying the annotation pipeline.

For the FACET dataset~\cite{gustafson2023facet}, an additional filtering step was applied: only samples for which the dataset's internal label-consensus score was at least 2.0 were retained. Labels below this threshold reflect disagreement among annotators and were excluded to avoid introducing noise during model evaluation.

% -----------------------------------------------------------
\subsection{Gender Annotation}
\label{sec:gender_annotation}

Five gender classifiers were evaluated on the pre-processed images:
\begin{enumerate}
	\item \textbf{HuggingFace Gender Classifier}~\cite{rizvandwiki2023gender} — a ViT-based binary classifier fine-tuned for gender prediction.
	\item \textbf{Realistic Gender Classifier}~\cite{prithivMLmods2024realistic} — a ViT-based model trained on high-quality portrait imagery.
	\item \textbf{DeepFace}~\cite{serengil2020deepface} — a facial analysis library offering gender prediction as one of several demographic attributes.
	\item \textbf{InsightFace}~\cite{deng2019insightface} — a face analysis toolkit with an integrated gender attribute predictor.
	\item \textbf{FairFace}~\cite{karkkainen2021fairface} — a model explicitly trained on a demographically balanced dataset to improve cross-group generalisation.
\end{enumerate}

All five models were evaluated on the Casual Conversations v2 (CCv2)~\cite{hazirbas2022ccv2} and FACET~\cite{gustafson2023facet} datasets after applying the pre-processing described above. The Realistic Gender Classifier achieved the highest overall performance on the CCv2 face-only subset, obtaining an overall accuracy of 88.6\%, with a male accuracy of 92.4\% and a female accuracy of 86.2\%. Per-skin-tone analysis revealed a performance gradient, with accuracy ranging from 75.6\% at MST~1 to 90.6\% at MST~5, declining again to 72.6\% at MST~10; female classification was consistently weaker than male classification across all skin tone groups. Despite this variation, the Realistic Gender Classifier outperformed all other candidates in both overall accuracy and consistency across the two evaluation datasets, and was therefore selected for final annotation.

% -----------------------------------------------------------
\subsection{Skin Tone Annotation}
\label{sec:skintone_annotation}

\paragraph{Scale selection.}
The first design decision concerned the choice of skin tone scale. The two most common options in the fairness literature are the Fitzpatrick scale~\cite{fitzpatrick1988skin}, which groups skin into six phototype categories, and the Monk Skin Tone (MST) scale~\cite{monk2023monk}, a ten-point scale developed with representational equity in mind. Following the recommendation of \cite{schumann2023consensus}, who found that the MST scale was perceived as more inclusive by darker-skinned respondents and by non-White and female survey participants, the MST scale was adopted for this work. Skin tone predictions were produced both as raw MST scores (1--10) and as coarsened three-bin groupings — Light (MST~1--3), Mid (MST~4--7), and Dark (MST~8--10) — following the binning convention of \cite{wu2023skintone}.

\paragraph{Model evaluation.}
Four skin tone estimation approaches were evaluated on the CCv2, FACET, and Monk Skin Tone Estimation (MSTE)~\cite{wu2023skintone} datasets:

\begin{enumerate}
	\item \textbf{SkinTone Classifier Library}~\cite{skintone2023lib} — an off-the-shelf Python library that uses colour histogram analysis to assign MST labels; applied directly without training.
	\item \textbf{RandomForest}~\cite{wu2023skintone} — a histogram-based classifier trained on colour features extracted from the segmented face region, following the methodology described by \cite{wu2023skintone}.
	\item \textbf{DenseNet-121}~\cite{wu2023skintone} — a convolutional network trained from scratch on the CCv2 dataset, following the architecture evaluated in \cite{wu2023skintone}.
	\item \textbf{VGG-16}~\cite{simonyan2015vgg} — a deep convolutional network fine-tuned for MST prediction, adapted from the architecture shown to perform well under diverse lighting conditions by \cite{krishnapriya2024lighting}.
\end{enumerate}

All models except the SkinTone Classifier Library were trained on the CCv2 dataset and evaluated on a held-out set of unseen images, as well as on the FACET and MSTE datasets to assess generalisation.

\paragraph{VGG-16 optimisation.}
VGG-16 achieved the best overall accuracy and generalisation across the three evaluation datasets. To further optimise the model, four binary design choices were explored, yielding eight model variants in total:

\begin{itemize}
	\item \textbf{Output formulation}: classification (cross-entropy loss over MST classes) versus regression (mean squared error over a continuous MST score clamped to $[1, 10]$).
	\item \textbf{Colour space}: RGB versus CIE L*a*b* (LAB), the latter having been shown to better separate luminance from chrominance and to improve robustness to lighting variation~\cite{krishnapriya2024lighting}.
	\item \textbf{Label granularity}: ten-class MST (MST10) versus three-class binned labels (MST3).
	\item \textbf{Background handling}: uniform black background versus an average skin-tone background computed from the unmasked pixels of each image.
\end{itemize}

Although the regression variants yielded lower exact-match accuracy than the corresponding classification variants, they produced substantially higher within-one accuracy (i.e.\ the proportion of predictions within one MST step of the ground truth) and higher three-bin accuracy. Since the downstream analysis operates on the coarsened three-bin groupings, and since regression provides a continuous score that degrades more gracefully near bin boundaries, the regression formulation was preferred. Among the remaining axes, the RGB colour space with a black background generalised more consistently across all three evaluation datasets than the LAB variants or the average-background alternatives. The final selected model was therefore the \textbf{VGG-16 MST10 Regression RGB Black-Background} variant.

\paragraph{Final annotation.}
The selected gender classifier and skin tone model were integrated into a unified annotation pipeline implemented in \texttt{GenderSkinToneAnnotation.py}. The pipeline applies MediaPipe FaceMesh segmentation to each detected face, runs the Realistic Gender Classifier on the RGB face crop and the VGG-16 MST regression model on the same crop, and writes the resulting gender label, continuous MST score, and three-bin skin tone category to a JSONL manifest. This pipeline was applied to the full sets of face-containing images retrieved from Re-LAION-5B and generated by Stable Diffusion v1.5, producing the demographic labels used in all subsequent analyses.